\newpage

\begin{center}
    {\Large\itshape Alban en sa « légende dorée »}
\end{center}

\vspace{1.2em}

\lettrine[lines=3,lhang=0.1,nindent=0em]{M}{ais} c'est un fait, le Seigneur de Nougaret, baron d'Alban, n'aurait pas passé une journée sans invoquer Notre Dame d'Alban -- « débonnaire Dame, plus savoureuse que le miel, plus suave que le parfum des violettes ». Ainsi s'exprimait le langage du temps\ldots Et ce dévot serviteur de Notre Dame n'en continuait pas moins ses vilenies. Hélas, les méfaits d'Alain devinrent si nombreux que bientôt, la révolte gronda dans la baronnie. Les paysans soulevés contre lui l'enfermèrent dans la prison de son propre château. Tête basse et l'œil triste, notre baron s'attendait à être pendu d'un moment à l'autre par ses propres sujets irrités. Mais Satan survint. Il promit la liberté. Une seule condition était exigée : livrer son âme au démon pour l'éternité. Le contrat tout prêt était à portée de la main. On n'avait qu'à signer et l'affaire était réglée. Troublante alternative que celle d'une mort imminente ou celle d'une liberté si chèrement achetée. Les Anglais prétendent que si l'on soupe avec le diable, il faut prendre une longue cuiller. Autant dire que discuter avec Satan, c'est être perdu.
